\subsubsection{Coverage}

This section presents the coverage planning algorithms employed in this study. While a comprehensive technical analysis of these algorithms is beyond the scope of this work, \newline
they are described with sufficient detail to enable understanding of their application within our framework. \newline
For in-depth technical details and formal derivations, the reader is referred to the original publications cited herein.

In order to demonstrate the different navigators, \newline
we use a map provided by the \href{https://github.com/MartijnWisse/mirte\_navigation}{mirte navigation} ros2 package. \newline
\texttt{Figure \{number\} \textless fig-original\_costmap\textgreater } shows the costmap generated by Nav2.

\begin{figure}[!htbp]
\centering
\includegraphics[width=0.7\linewidth]{figures/original_costmap.*}
\end{figure}

\begin{verbatim}
- - - - - - - - - - - - - - - - - - - - - - - - - - - - - - - - - - - - - - - - - - - - - - - - - - - - - - - - - - - - - - - - - - - - - - - - - - -
ModuleNotFoundError                       Traceback (most recent call last)
Cell In[6], line 10
      7 from matplotlib.widgets import Slider
      9 from local_codebase.imports.navigator_types import SystematicNavigator
 - - -> 10 from local_codebase.imports.navigators import BousPath, SkeletonPath
     13 def show_imgs(imgs, labels):
     14     n = len(imgs)

File c:\Users\maxva\Desktop\Bep Eindproject\Lab -Cleanup -Robot -using -the -Mirte -Master -Platform\content\local_codebase\imports\navigators.py:8
      5 import local_codebase.imports.utils as ut
      6 from local_codebase.imports.utils import LogType
 - - - -> 8 from trajgenpy import Geometries
     10 import shapely
     11 from shapely.validation import explain_validity

ModuleNotFoundError: No module named 'trajgenpy'
\end{verbatim}

\paragraph{The Free space Polygon}

Each defined navigator must update its internal map and as a result plan a path over this map.

A functino ``\texttt{update\_map()}'' first takes the costmap and, given a threshold, sets the free space to 0 and occupied space to 1. \newline
It then detects the boundaries of the free space and saves these as contours.

\begin{verbatim}
- - - - - - - - - - - - - - - - - - - - - - - - - - - - - - - - - - - - - - - - - - - - - - - - - - - - - - - - - - - - - - - - - - - - - - - - - - -
NameError                                 Traceback (most recent call last)
Cell In[7], line 3
      1 navigator = SystematicNavigator()
 - - - -> 3 navigator.update_map(map_img)
      4 contours = navigator.contours
      6 color_img = cv2.cvtColor(map_img, cv2.COLOR_GRAY2RGB)

NameError: name 'map_img' is not defined
\end{verbatim}

\paragraph{Boustrophedon Coverage}

The first navigator that will be covered is a Boustropedon based Coverage planner. \newline
This navigator plans in two steps:

\begin{itemize}
\item Decomposition
\item Swath Generation
\end{itemize}

\subparagraph{Boustrophedon Cellular Decomposition}

To effectively cover the entire map with boustrophedon (back and forth) swaths, the map is first partitioned (decomposed) into smalller areas that can be covered fully with these swaths. This is done by sweeping a straight line across the map and creating new cells where the connectivity of this line decreases. Cells are closed and a new one opens when this connectivity increases again or the line is fully restored, The results of Boustrophedon cellular Decomposition can be viewed in the cell below.

\begin{verbatim}
- - - - - - - - - - - - - - - - - - - - - - - - - - - - - - - - - - - - - - - - - - - - - - - - - - - - - - - - - - - - - - - - - - - - - - - - - - -
NameError                                 Traceback (most recent call last)
Cell In[11], line 2
      1 # Update internal map
 - - - -> 2 navigator = BousPath()
      3 navigator.update_map(map_img)
      4 
      5 # Perform Boustropethan Celular Decomposition

NameError: name 'BousPath' is not defined
\end{verbatim}

\subparagraph{Boustrophedon Path Generation}

After decomposition, each cell is populated with a boustrophedon path. Swaths are generated as straight lines parallel to the line used during decomposition. The transitions, often called headlands, are then generated to create one continuous path inside the cell.

\begin{verbatim}
Number of cells: 16
generating full trajectory
\end{verbatim}

\includegraphics[width=0.7\linewidth]{files/2d7b2cbc1b8dc1aa7aa5a49495009de2.png}

\paragraph{SkeletonPath}

The \texttt{SkeletonPath} planner is a topology-based coverage planner that generates navigation paths using the morphological skeleton of the environment.

Unlike the \texttt{BousPath} planner, which decomposes the free space into sweepable cells, the skeleton planner reduces the free space to a graph-like centerline representation. This representation preserves the connectivity of the environment while simplifying path planning.

The implementation is inspired by recent work in morphology-based exploration and autonomous navigation.


\bigskip
\centerline{\rule{13cm}{0.4pt}}
\bigskip

\subparagraph{Planner Initialization}

The planner can be initialized similarly to the other navigators:

\subparagraph{Reading and Skeletonizing the Map}

The skeleton planner generates a morphological skeleton.

Internally, the planner fills the free-space contours, skeletonizes the free-space region, extracts skeleton pixels,
and converts them into Cartesian waypoints. The generated waypoints are then stored.

\begin{verbatim}
- - - - - - - - - - - - - - - - - - - - - - - - - - - - - - - - - - - - - - - - - - - - - - - - - - - - - - - - - - - - - - - - - - - - - - - - - - -
NameError                                 Traceback (most recent call last)
Cell In[3], line 1
 - - - -> 1 navigator.read()
      3 skeleton_img = cv2.addWeighted(
      4     navigator.contour_img,
      5     0.8,
   (...)
      8     0,
      9 )
     11 # Convert grayscale map to RGB so we can draw colored points

NameError: name 'navigator' is not defined
\end{verbatim}

\subparagraph{Path Planning}

After the waypoints have been generated, they are then sampled and ordered. \newline
The algorithm does this by first constructing a graph representation of the skeleton tree. \newline
The waypoints then become graph nodes. Edges are generated between the nodes with a certain distance threshold {\textbackslash}

To efficiently identify neighboring points, a k-d tree is used for spatial indexing. \newline
Starting from the closest waypoint to the robot, the algorithm recursively explores nearby nodes and connects them within the graph.

Leaf nodes are then identified in the tree, (nodes with one connection).

\includegraphics[width=0.7\linewidth]{files/4599d3c4df780082d726948f94d0cc2d.png}

\subparagraph{Path Ordering}

First the closest leaf node to the robot is found, this is the starting position. \newline
The planner traces a path along the tree and selects the shortest one to the next leaf node. \newline
It then appends these points to the path and sets this new leaf node as the current one and marks it as visited. \newline
The algorithm repeats this process untill all leaf nodes have been explored successfully.

\begin{verbatim}
- - - - - - - - - - - - - - - - - - - - - - - - - - - - - - - - - - - - - - - - - - - - - - - - - - - - - - - - - - - - - - - - - - - - - - - - - - -
NameError                                 Traceback (most recent call last)
Cell In[5], line 5
      1 # - - - - - - - - - - - - - - - - - - - - - - - - - - - - - - - - - - - - - - - - - - - - - - - - - -
      2 # Generate planner path
      3 # - - - - - - - - - - - - - - - - - - - - - - - - - - - - - - - - - - - - - - - - - - - - - - - - - -
      4 
 - - - -> 5 navigator.start = [0, 0]
      6 navigator.plan_path()
      7 
      8 path_points = navigator.path

NameError: name 'navigator' is not defined
\end{verbatim}